\bookchapter{Bab 57: Jalan Menuju Tingkat Lanjut — Izin, Notifikasi, dan Arsitektur Modern}

\chapterintro{Bisa membuat App hanyalah awal. Bab ini membahas beberapa rintangan dari “bisa membuat” menuju “membuat dengan baik”: cara meminta izin dengan elegan, cara mengirim notifikasi yang tidak menyebalkan, cara menaikkan arsitektur satu tingkat, dan — di tahun 2026 — cara menulis kode Android berpasangan dengan AI. Semuanya akan segera kamu jumpai di proyek nyata.}

\booksection{Izin saat runtime: Ketuk pintu dulu, baru masuk}

Kamera, lokasi, kontak, notifikasi... “zona sensitif” ini wajib diminta App kepada pengguna saat berjalan (sejak Android 6). Aturannya: deklarasikan dulu (Manifest), lalu minta (di kode), dan bila pengguna menolak, turunkan fitur dengan elegan. Di Compose, yang paling praktis memakai rememberPermissionState dari pustaka resmi Accompanist:

\begin{lstlisting}[language=Kotlin]
// Deklarasikan dulu di Manifest:
// <uses-permission android:name="android.permission.CAMERA" />

@Composable
fun CameraButton() {
    val permission = rememberPermissionState(
        permission = android.Manifest.permission.CAMERA
    )
    Button(onClick = {
        when {
            permission.status.isGranted -> openCamera()
            permission.status.shouldShowRationale -> showExplainDialog()
            else -> permission.launchPermissionRequest()
        }
    }) { Text("Jepret") }
}
\end{lstlisting}

\begin{tcbwarn}{Perhatian: Bagaimana kalau ditolak permanen?}
Setelah pengguna mencentang “jangan tanya lagi”, launchPermissionRequest tidak akan memunculkan dialog lagi. Saat itu, arahkan pengguna menyalakannya manual di Pengaturan sistem (lompat ke halaman pengaturan aplikasi dengan Intent) dan jelaskan di antarmuka “izin ini dipakai untuk apa”. Sikap yang ramah membuat pengguna jauh lebih bersedia memberi izin — teks permintaan izin layak ditulis dengan serius.

\end{tcbwarn}

\booksection{Notifikasi: Mengucapkan hal berguna di waktu yang tepat}

Notifikasi adalah saluran App menjangkau pengguna, sekaligus fitur yang paling gampang “diblokir” pengguna. Tiga aturan mengirim notifikasi: hanya kirim yang berguna, biarkan pengguna bisa mematikan, jangan ganggu tengah malam. Pemakaian dasarnya:

\begin{lstlisting}[language=Kotlin]
// Android 13+ wajib minta izin POST_NOTIFICATIONS dulu (lihat bagian sebelumnya)

// 1. Buat channel notifikasi (wajib di Android 8+)
val channel = NotificationChannel(
    "reminder", "Pengingat", NotificationManager.IMPORTANCE_DEFAULT
)
val manager = getSystemService(NotificationManager::class.java)
manager.createNotificationChannel(channel)

// 2. Kirim notifikasi
val notification = NotificationCompat.Builder(this, "reminder")
    .setSmallIcon(R.drawable.ic_notification)
    .setContentTitle("Pengingat mencatat")
    .setContentText("Hari ini belum mencatat, lho")
    .setPriority(NotificationCompat.PRIORITY_DEFAULT)
    .build()
if (ActivityCompat.checkSelfPermission(this,
        Manifest.permission.POST_NOTIFICATIONS) == PackageManager.PERMISSION_GRANTED
) {
    manager.notify(1, notification)
}
\end{lstlisting}

Tugas terjadwal (misalnya “pengingat mencatat pukul 20.00 tiap hari”) gunakan WorkManager — ia “asisten penjadwal latar belakang” yang direkomendasikan resmi, menjamin tugas dieksekusi di waktu yang tepat tanpa kamu repot mengurus optimasi baterai. Untuk kebutuhan sekelas alarm, memakai ini pasti tepat.

\booksection{Arsitektur lanjutan: Dependency injection dan modularisasi}

Begitu proyek membesar, dua “masalah manis” akan mendatangi kamu:

\begin{itemize}
  \item Pembuatan objek kacau: tiap halaman membuat database dan API sendiri-sendiri — kelola semuanya dengan Hilt (framework dependency injection); "siapa butuh apa, tinggal minta, langsung disodorkan";
  \item Kode menumpuk jadi gunung: proyek satu modul makin lambat tiap kali diubah — pecah menjadi multi-modul (app / core / feature-xxx), tiap fitur dikembangkan dan diuji secara mandiri.
\end{itemize}

\begin{tcbinfo}{Tips: Jangan over-engineering}
Hilt dan modularisasi memang alat yang bagus, tetapi untuk proyek satu orang jangan buru-buru memakainya. Aplikasi catatan keuangan dengan satu modul + pembuatan dependensi manual sudah sangat cukup. Tunggu sampai kode melewati puluhan ribu baris, kompilasi mulai melambat, atau mulai kerja sama beberapa orang, baru masukkan sesuai kebutuhan. Alat melayani tujuan, bukan untuk dipamerkan — ini juga prinsip “kesederhanaan di atas segalanya” di buku pemrograman AI kita.

\end{tcbinfo}

\booksection{Cara bermain baru 2026: Pengembangan Android berbantuan AI}

Di tahun 2026, AI sudah menyatu dalam alur pengembangan Android. Kalau ingin mempercepat, rangkaian kombinasi ini sudah sangat matang:

\begin{center}
\begin{tabularx}{\linewidth}{|X|X|X|}
\hline
\textbf{Tahap} & \textbf{Cara AI membantu} & \textbf{Tugasmu} \tabularnewline
\hline
Menulis kode & Gemini bawaan Android Studio / plugin AI pihak ketiga untuk pelengkapan dan pembuatan kode & Menjelaskan kebutuhan dengan jelas, meninjau hasilnya \tabularnewline
Mencari referensi & Melemparkan dokumentasi resmi dan pesan error kepada AI untuk dijelaskan & Memverifikasi kesimpulan penting kembali ke dokumentasi resmi \tabularnewline
Menulis test & Meminta AI membuat test case menurut tiga kategori "normal / batas / tak terduga" & Memastikan test case sesuai makna bisnisnya \tabularnewline
Memperbaiki bug & Menempelkan stacktrace + kode terkait, AI menemukan lokasi dan memperbaiki & Memahami perubahan, menjalankan test untuk verifikasi \tabularnewline
\hline
\end{tabularx}
\end{center}

\begin{center}\small\color{gray}Empat skenario pengembangan Android berbantuan AI\end{center}

Metodologi konkretnya — cara memberi instruksi ke AI, cara menilai hasil, cara menghindari jebakan — dibahas lengkap dalam 11 bab di Bagian 7 buku ini, “Praktik Pemrograman AI” (Bab 37-47). Sangat disarankan dibaca beriringan: sintaks Android diserahkan ke bagian ini, kiat kolaborasi dengan AI diserahkan ke Bagian 7.

\booksection{Optimasi performa dan alat debugging Compose}

Menulis dengan Compose memang asyik, tetapi “rekomposisi berlebihan” juga bisa membuat tersendat. Tiga trik praktis:

\begin{itemize}
  \item Stabilkan parameter: tambahkan @Immutable pada data class atau gunakan derivedStateOf agar perubahan state yang tak berhubungan tidak memicu rekomposisi luas;
  \item Ingat pakai key: items(list, key = \{ it.id \}) membuat daftar hanya memperbarui item yang memang harus berubah saat ada penambahan/penghapusan — tanpa key, menghapus satu item bisa menggambar ulang seluruh daftar;
  \item Mode Compose di Layout Inspector: melihat jumlah rekomposisi tiap Composable secara waktu nyata; siapa yang sering rekomposisi terlihat jelas, dan optimasi tinggal mengikuti titik panasnya.
\end{itemize}

Selain itu, aktifkan laporan compiler Compose (Metrics) untuk menghasilkan laporan “komponen mana yang tidak stabil” — di jalan menuju tingkat lanjut, data selalu lebih andal daripada perasaan.

\begin{tcbpractice}{Praktik: Tambahkan "paket tiga serangkai" ke aplikasi catatan keuangan}
(1) Tambahkan “pengingat mencatat pukul 20.00 tiap hari” ke aplikasi catatan keuangan (WorkManager + saluran notifikasi + izin notifikasi Android 13); (2) ubah pembuatan database dan API menjadi injeksi Hilt; (3) coba percepat implementasi fitur “filter bulanan” dengan alat AI, lalu periksa hasilnya dengan metode penilaian tiga lapis. Setelah tiga serangkai ini selesai, proyekmu sudah tampak seperti “pasukan reguler”.

\end{tcbpractice}

\begin{tcbsummary}{Ringkasan bab}
\begin{itemize}
  \item Izin sensitif: deklarasi di Manifest + permintaan saat runtime + arahan elegan setelah ditolak
  \item Tiga aturan notifikasi: hanya yang berguna, bisa dimatikan, tidak mengganggu; saluran + izin dua-duanya wajib ada
  \item Tugas terjadwal pakai WorkManager, jangan bikin roda sendiri
  \item Hilt mengurus objek, modularisasi mengurus kode — proyek kecil satu orang jangan buru-buru
  \item Empat skenario bantuan AI: menulis kode, mencari referensi, menulis test, memperbaiki bug; metodologinya di Bagian 7 (Bab 37-47)
\end{itemize}

\end{tcbsummary}
